% Options for packages loaded elsewhere
\PassOptionsToPackage{unicode}{hyperref}
\PassOptionsToPackage{hyphens}{url}
%
\documentclass[
]{article}
\usepackage{amsmath,amssymb}
\usepackage{lmodern}
\usepackage{iftex}
\ifPDFTeX
  \usepackage[T1]{fontenc}
  \usepackage[utf8]{inputenc}
  \usepackage{textcomp} % provide euro and other symbols
\else % if luatex or xetex
  \usepackage{unicode-math}
  \defaultfontfeatures{Scale=MatchLowercase}
  \defaultfontfeatures[\rmfamily]{Ligatures=TeX,Scale=1}
\fi
% Use upquote if available, for straight quotes in verbatim environments
\IfFileExists{upquote.sty}{\usepackage{upquote}}{}
\IfFileExists{microtype.sty}{% use microtype if available
  \usepackage[]{microtype}
  \UseMicrotypeSet[protrusion]{basicmath} % disable protrusion for tt fonts
}{}
\makeatletter
\@ifundefined{KOMAClassName}{% if non-KOMA class
  \IfFileExists{parskip.sty}{%
    \usepackage{parskip}
  }{% else
    \setlength{\parindent}{0pt}
    \setlength{\parskip}{6pt plus 2pt minus 1pt}}
}{% if KOMA class
  \KOMAoptions{parskip=half}}
\makeatother
\usepackage{xcolor}
\setlength{\emergencystretch}{3em} % prevent overfull lines
\providecommand{\tightlist}{%
  \setlength{\itemsep}{0pt}\setlength{\parskip}{0pt}}
\setcounter{secnumdepth}{-\maxdimen} % remove section numbering
\ifLuaTeX
  \usepackage{selnolig}  % disable illegal ligatures
\fi
\IfFileExists{bookmark.sty}{\usepackage{bookmark}}{\usepackage{hyperref}}
\IfFileExists{xurl.sty}{\usepackage{xurl}}{} % add URL line breaks if available
\urlstyle{same} % disable monospaced font for URLs
\hypersetup{
  hidelinks,
  pdfcreator={LaTeX via pandoc}}

\author{Dietmar Gerald Schrausser}
\date{09.05.2026}


\begin{document}

\hypertarget{foliumblat-iiv}{%
\section{Foliu{[}m{]}/Blat IIv}\label{foliumblat-iiv}}

Eine zeitgenössische deutsche Übersetzung wird präsentiert, die auf
einem Vergleich der bestehenden lateinischen, frühneuhochdeutschen
(ENHG) und englischen (Übersetzungs-) Texte basiert, um ein besseres
Verständnis zu ermöglichen. Insbesondere da der ENHG-Text im Vergleich
zum modernen Hochdeutschen schwer verständlich ist, da dieser fundierte
Deutschkenntnisse (als Muttersprache) sowie Kenntnisse verschiedener
Dialekte voraussetzt.

\hypertarget{vom-werk-des-ersten-tages}{%
\subsection{Vom Werk des ersten Tages}\label{vom-werk-des-ersten-tages}}

\begin{quote}
\textbf{D}Ixit deus. fiat lux.\\
Et facta e{[}st{]} lux. Et vidit deus lucem q{[}uia{]} esset bona.\\
Et diuisit luce{[}m{]} a tenebris.\\
Appelauitq{[}ue{]} luce{[}m{]} die{[}m{]}.\\
Et tenebras nocte{[}m{]}.\\
Factu{[}m{]}q{[}ue{]} est vespere {[}et{]} mane dies vnus.
\end{quote}

Mit göttlichem, übermenschlichen Fleiß verfasste \textbf{Moses} ein
wunderbares Meisterwerk über alle Geheimnisse der Natur, ein Buch das
alle anderen an Eloquenz und Genialität\footnote{`eloq{[}ue{]}ntia{[}m{]}
  {[}et{]} ingeniu{[}m{]}' und `außprechlichkeit vn{[}d{]}
  sinnreichikeit'.} übertrifft: Denn der glorreiche Gott, der das
\emph{wahre Licht ist}, das Licht wertschätzend, alle Dinge im Licht
machend, begann die Erschaffung der \emph{Welt}
korrektesterweise\footnote{`rectissime' und `gar recht'.} aus dem
Licht\footnote{`a luce' und `am liecht'.}.\footnote{wird Kaiser Karl dem
  Großen zugeschrieben, sein Reformprogramm `Admonitio Generalis' (s.
  Baluzius,
  \href{https://digital.ub.uni-paderborn.de/eab/content/zoom/1075337}{1677})
  bezieht sich ausschließlich auf die Bibel und das kanonische Recht und
  spiegelt die Grundlagen der \emph{karolingischen Renaissance} wider.}
Diese vollendete einen natürlichen Tag in ihrem Kreislauf\footnote{`circuito'
  und `umkreis'.} über drei Tage bis zum vierten (in dem die
Leuchten\footnote{`lu{[}m{]}inaria formata' und `großen liecht geformt'.}
geformt sind).\footnote{von den `Etymologiae' (De Seville,
  \href{https://books.google.com/books?id=wb8Z_vlSyJwC}{1802}, Buch V,
  Kapitel 30).}

Unter allen physischen Dingen das edelste und vom spirituellen Wesen das
Naheliegenste und Optimum, in seiner Schönheit\footnote{`pulcritudinem'
  und `das sein schone'.} allumfassend\footnote{`maxi{[}m{]}e
  co{[}mun{]}icans' und `allermeist gemainsam'.}, selbst den kleinsten
Punkt der Welt erfüllend, ist es ausschließlich das \emph{Licht},
welches das Unreinen schadlos durchquert, wodurch die ganze Welt gut und
schön ist. Verdientermaßen sah er das Licht das gut war, es sei nichts
anderes, als ein Abbild, feines und schattiges Gleichnis des obersten
Guten.\footnote{Passage aus dem `Heptaplus' (Mirandola \& Mirandola,
  \href{https://doi.org/10.3931/e-rara-61217}{1601}, Buch I, Kapitel 4).}

Sobald nun der Geist diese Wasser erreichte, den Unterbau\footnote{`s{[}u{]}biectu{[}m{]}'
  und `underwurff'.} durchdrang, entstand auf Befehl des
Architekten\footnote{`artificis' und `werckmeisters'.} das Licht in
Schönheit und Pracht, es erstrahlte wie Lichtwolken, welche die oberen
Regionen der Welt mit ihrem Glanz erleuchten (wie die heranbrechende
Morgensonne), den oberen, dann in Folge ebenso den unteren Teil der
\emph{Sphäre} erhellend.\footnote{aus dem `Heptaplus' (Mirandola \&
  Mirandola, \href{https://doi.org/10.3931/e-rara-61217}{1601}, Buch 1,
  Kapitel 2).} Danach trennte er diese, sodass Finsternis und Licht zwei
unterschiedliche Hemisphären hätten. Wegen seiner Klarheit, welche die
Finsternis reinigt\footnote{`purgat' und `rainigt'.} nannte er das Licht
den Tag\footnote{`Appellauit luce{[}m{]} die{[}m{]}' vgl. `Dimmen' und
  `liecht hies er den tag'.}, die schädliche Finsternis hingegen
Nacht\footnote{`tenebras a noce{[}n{]}do nocte{[}m{]}' und `finsternus
  vo{[}n{]} beschedigung'.}, damit sie die Augen nicht
sähen\footnote{Rein physikalisch betrachtet erscheint dies unlogisch,
  insbesondere da übermäßig starkes Licht wohl schädliche Auswirkungen
  hätte, was bei übermäßiger Dunkelheit kaum der Fall wäre. Im
  übertragenen Sinne hingegen erscheint es nachvollziebar, da
  \emph{Licht} im Sinne von \emph{Wissen} eine \emph{Dunkelheit} im
  Sinne von \emph{Unwissenheit} durchaus \emph{erhellt} und somit
  \emph{reinigt}.}.\footnote{aus der `Historia Scholastica' (Comestor,
  \href{https://books.google.com/books?id=fKkLGynZVoYC}{1699}, Liber
  Genesis, Kapitel 3).} Aus Dimensionierung dieser Teile schuf er Tag
und Nacht, wobei die abwechselnde Folge der Räume und Kreislaufe
\emph{die} ewigen, sich wiederholenden Zeiten bilden, die wir Jahre
nennen.\footnote{aus den `Divinae Institutiones' (Lactantius,
  \href{https://books.google.com/books?id=qGsnghsi3uUC}{1497}, Buch II,
  Kapitel 9 oder 10).}

So ist \emph{ein} Tag vergangen, der erste Tag der
\emph{Welt}\footnote{`seculi' (!).}, jedoch nicht der erste aller
Tage\footnote{`p{[}ri{]}m{[}us{]} dio{[}rum{]} o{[}mn{]}i{[}u{]}m' und
  `der erst aller tag'.}. Deshalb wird er nicht erster, sondern
\emph{ein} Tag genannt. Wahrhaftig\footnote{`sic'.}, an diesem Tag schuf
Gott die formlose\footnote{`informe{[}m{]}'.} Materie, die Engel, den
Himmel, das Licht, die Erde, das Wasser, die Luft (usw.). Weiters wurden
zwei einander entgegengesetzte, unterschiedliche Teile der Erde, der
Osten und der Westen geschaffen.\footnote{aus den `Divinae
  Institutiones' (Lactantius,
  \href{https://books.google.com/books?id=qGsnghsi3uUC}{1497}, Buch II,
  Kapitel 9 oder 10).}

Der \emph{Osten} wird Gott zugeschrieben, er ist der
\emph{Quell}\footnote{`fons' und `prun{[}n{]}'.} des Lichts und der
Illustrator der Dinge, welcher uns zum ewigen Leben auferstehen lässt.
Der \emph{Westen} aber wird dem streitsüchtigen, entrüsteten und
boshaften Gemüt\footnote{`{[}con{]}turbate illi p{[}ra{]}ueq{[}ue{]}
  me{[}n{]}ti' und `zerstreitten entrüsten vnd boßhafftigen gemüt'.}
zugeschrieben, er bringt Finsternis, was darauf abzielt, die Menschen
durch die Sünde zu verführen und zu töten. Gleich wie das Licht aus dem
Osten entspringt und die weltliche Vernunft\footnote{`vite ratio' und
  `vernunfft des lebens'.} im Licht sich \emph{wendet}\footnote{`v{[}er{]}sat{[}ur{]}'
  und `swebt' (!).}, so kommen die Finsternisse aus dem Westen, dort
sind Tod und Zerstörung.\footnote{aus den `Divinae Institutiones'
  (Lactantius,
  \href{https://books.google.com/books?id=qGsnghsi3uUC}{1497}, Buch II,
  Kapitel 9 oder 10).}

Danach bestimmte Gott die anderen Regionen, den Süden und den Norden in
der selben Dimension, den ersten beiden Teilen entsprechend. So ist die
Region, die von der Sonne stark erwärmt wird, dem Osten am nächsten. Das
hingegen, was permanent gefroren und von Kälte betäubt ist entspricht
einem Teil des extremen Westens\footnote{`cui{[}us{]} extrem{[}us{]}
  occasus' und `des letzte{[}n{]} niderga{[}n{]}gs'.}. Wie Finsternis
dem Licht, so ist Kälte der Wärme entgegengesetzt. Wie die Wärme dem
Licht am nächsten ist, so der Süden dem Osten, sowie Kälte der
Finsternis am nächsten, so auch die Plage des Nordens\footnote{`plaga
  septe{[}m{]}trionis'.} dem Westen. All dies wird im \emph{Werk des
vierten Tages} verdeutlicht.\footnote{aus den `Divinae Institutiones'
  (Lactantius,
  \href{https://books.google.com/books?id=qGsnghsi3uUC}{1497}, Buch II,
  Kapitel 9 oder 10).}

\hypertarget{references}{%
\subsection{References}\label{references}}

Baluzius, S., ed.~(1677). Capitvlaria Karoli Magni. In \emph{Capitvlaria
Regvm Francorvm}, 1:189. Parisiis: Muguet.
\url{https://digital.ub.uni-paderborn.de/eab/content/zoom/1075337}

Comestor, P. (1699). \emph{Eruditissimi Viri Magistri Petri Comestoris
Historia Scholastica: Opus Eximium Magnam Sacrae Scripturae Partem, Quae
{[}Et{]} in Serie, {[}Et{]} in Glossis Crebro Diffusa Erat, Breviter
Complectens}. Matriti: ex officina Antonij Gonçalez de Reyes \ldots{} :
a costa de Francisco Sacédon \ldots{}
\url{https://books.google.com/books?id=fKkLGynZVoYC}

De Seville, I. (1802). \emph{S. Isidori Hispalensis Episcopi Hispaniarum
Doctoris Opera Omnia}. Edited by Arevalo, F. Romae: Apvd Antonivm
Fvlgonivm. \url{https://books.google.com/books?id=wb8Z_vlSyJwC}

Lactantius, L.C.F. (1497). \emph{Divinae Institutiones}. Venetiis: Simon
Bevilaqua. \url{https://books.google.com/books?id=qGsnghsi3uUC}

Mirandola, G., \& Mirandola, G. F. (1601). \emph{Ioannis Pici,
Mirandulae \ldots{} opera quae extant omnia \ldots{}} Basileae: per
Sebastianum Henricpetri. \url{https://doi.org/10.3931/e-rara-61217}

\end{document}
